% ===================================================================
%  calibrate-io.tex — TOUTES les mesures du fac-simile « Expliko-Libreto di la Delmas-Tabeli helpanta ».
%  Fichier PRODUIT par outils/kalibro.py : ne pas le modifier a la main.
%  Une seule constante y est physique — la largeur du papier — et elle
%  est passee en argument. Tout le reste vient de outils/inv-io.json.
%
%  Releve : 62 pages de 35 lignes ou plus.
%    justification  1456 px  (ecart-type 320)
%    hauteur bloc   2762 px
%    pas des lignes 75.25 px  (ecart-type 15.41)
%    image de page  1917 x 3300 px
%    marge gauche   206 px
%    1re ligne a    120 px du bord
%  Echelle retenue : 18.333 px/mm, prise sur une hauteur de papier
%  de 180 mm (notice de bibliotheque) ; la page mesure alors
%  104.6 mm de large.
% ===================================================================

\newcommand{\VUlangue}{io}
\newcommand{\VUmarque}{Gilles-Philippe Morin kompris, skanis e direktis la transskribo di ca libro.}

% 1. GEOMETRIE
\newcommand{\VUpapierLargeur}{104.56mm}
\newcommand{\VUpapierHauteur}{180.00mm}
\newcommand{\VUtexteLargeur}{79.42mm}
\newcommand{\VUtexteHauteur}{150.65mm}
% LE FOLIO EST LA PREMIERE ENCRE DE LA PAGE, LE TEXTE VIENT APRES.
% Premiere version : les deux ordonnees etaient egales, et le folio se
% composait DANS la premiere ligne de texte — « 6 » barrait « esas » au
% folio 6. La mesure de outils/inventaire.py donne le haut du BLOC
% D'ENCRE, c'est-a-dire le folio sur toute page foliotee ; le corps du
% texte commence un rang plus bas. On ajoute donc un pas de ligne, et
% les pages sans folio (ouvertures de tableau) commencent a la meme
% ordonnee que les autres, comme au fac-simile.
\newcommand{\VUmargeSup}{10.68mm}
\newcommand{\VUfolioY}{6.57mm}

% 2. CORPS ET INTERLIGNAGE
\newcommand{\VUinterligne}{11.68pt}
\newcommand{\VUcorps}{10.93pt}
\newcommand{\VUcorpsNote}{9.09pt}
\newcommand{\VUinterligneNote}{9.72pt}
% Renfoncement d'alinea : provisoire, a relever (outils/mesures.py).
\newcommand{\VUalinea}{3.97mm}
\newcommand{\VUalineaNote}{3.18mm}
\newcommand{\VUblancAlineaValeur}{2.22pt}
\newcommand{\VUblancNote}{2.92pt}
\newcommand{\VUfiletnoteLargeur}{21.44mm}

% 3. ESPACE-MOT
% Largeur proche de celle de la fonte, elasticite large : chaque ligne
% trouve sa mesure sans que TeX ait a couper ailleurs qu'au point releve.
\newcommand{\VUespaceRatio}{0.32}
\newcommand{\VUespaceEtire}{0.30}
\newcommand{\VUespaceSerre}{0.14}

% 4. FONTES
\newcommand{\VUfonte}{XCharter-TLF}
\newcommand{\VUfonteGras}{ntxtlf}
\newcommand{\VUratioGras}{0.98}
